\documentclass[CJK]{beamer}
\usepackage{CJKutf8}
\usepackage{beamerthemesplit}
\usetheme{Malmoe}
\useoutertheme[footline=authortitle]{miniframes}
\usepackage{amsmath}
\usepackage{amssymb}
\usepackage{graphicx}
\usepackage{color}
\usepackage{slashed}
\usepackage{simplewick}
\graphicspath{{../figures/}}
\def\be{\begin{equation}}
\def\ee{\nonumber\end{equation}}
\def\bea{\begin{eqnarray}}
\def\eea{\nonumber\end{eqnarray}}
\def\ii{{\dot{\imath}}}
\def\bch{\begin{CJK}{UTF8}{gbsn}}
\def\ech{\end{CJK}}
\def\bex{\begin{minipage}{0.3\textwidth}\includegraphics[width=1in]{jugelizi.png}\end{minipage}\begin{minipage}{0.6\textwidth}}
\def\eex{\end{minipage}}
\def\chtitle#1{\frametitle{\bch#1\ech}}
\def\skipline{{\vskip0.1in}}
\def\skiplines{{\vskip0.2in}}
\def\lagr{{\mathcal{L}}}
\def\hamil{{\mathcal{H}}}
\def\vecv{{\mathbf{v}}}
\def\vecx{{\mathbf{x}}}
\def\veck{{\mathbf{k}}}
\def\vecp{{\mathbf{p}}}
\def\vecn{{\mathbf{n}}}
\def\vecA{{\mathbf{A}}}
\def\vecP{{\mathbf{P}}}
\def\vecsigma{{\mathbf{\sigma}}}
\def\hatJn{{\hat{J_\vecn}}}
\def\hatJx{{\hat{J_x}}}
\def\hatJy{{\hat{J_y}}}
\def\hatJz{{\hat{J_z}}}
\def\hatj#1{\hat{J_{#1}}}
\def\hatphi{{\hat{\phi}}}
\def\hatq{{\hat{q}}}
\def\hatpi{{\hat{\pi}}}
\def\vel{\upsilon}
\def\Dint{{\mathcal{D}}}
\def\adag{{\hat{a}^\dagger}}
\def\bdag{{\hat{b}^\dagger}}
\def\cdag{{\hat{c}^\dagger}}
\def\ddag{{\hat{d}^\dagger}}
\def\hata{{\hat{a}}}
\def\hatb{{\hat{b}}}
\def\hatc{{\hat{c}}}
\def\hatd{{\hat{d}}}
\def\hatN{{\hat{N}}}
\def\hatH{{\hat{H}}}
\def\hatp{{\hat{p}}}
\def\Fup{{F^{\mu\nu}}}
\def\Fdown{{F_{\mu\nu}}}
\def\newl{\nonumber \\}
\def\SIkm{\mathrm{km}}
\def\SIyr{\mathrm{yr}}
\def\SIGyr{\mathrm{Gyr}}
\def\SIeV{\mathrm{eV}}
\def\SIGeV{\mathrm{GeV}}
\def\SIm{\mathrm{m}}
\def\SIcm{\mathrm{cm}}
\def\SIJ{\mathrm{J}}
\def\SIs{\mathrm{s}}
\def\SIkg{\mathrm{kg}}
\def\SIg{\mathrm{g}}
\def\vece{\mathrm{e}}
\def\bmat#1{\left(\begin{array}{#1}}
\def\emat{\end{array}\right)}
\def\bcase#1{\left\{\begin{array}{#1}}
\def\ecase{\end{array}\right.}
\def\calM{{\mathcal{M}}}
\def\calT{{\mathcal{T}}}
\def\calR{{\mathcal{R}}}
\def\barpsi{\bar{\psi}}
\def\baru{\bar{u}}
\def\barv{\bar{\upsilon}}
\def\bmini#1{\begin{minipage}{#1\textwidth}}
\def\emini{\end{minipage}}
\def\qeq{\stackrel{?}{=}}
\def\torder#1{\mathcal{T}\left(#1\right)}
\def\rorder#1{\mathcal{R}\left(#1\right)}


\title{Quantum Field Theory I \\ Lesson 12 - The First Feynman diagram}
\author{}
\date{}


\begin{document}

\begin{frame}
 
\begin{center}
\begin{Large}
\bch
量子场论 I 

{\vskip 0.3in}

第十二课 第一个Feynman图

\ech
\end{Large}
\end{center}

\vskip 0.2in

\bch
课件下载
\ech
https://github.com/zqhuang/SYSU\_QFTI

\end{frame}


\begin{frame} 
\chtitle{相互作用(Interaction)绘景} 
\bch
假设Hamilton算符可以拆成两部分: $\hat{H} = \hat{H}_0 + \hat{H}_I$, 则在各种绘景下态和可观测量算符的演化方程为：

\begin{tabular}{c|c|c|c}
\hline
\hline
            & Schr\"odinger & Heisenberg  & Interaction \\
\hline
State         & $\frac{d |\psi\rangle}{dt} = - \ii \hat{H}|\psi\rangle$ & constant & $\frac{d |\psi\rangle}{dt} = - \ii \hat{H}_I|\psi\rangle$  \\
\hline
Observable & const. & $\frac{d\hat{O}}{dt} = -\ii [ \hat{O}, \hat{H} ]$ & $\frac{d\hat{O}}{dt} = -\ii [ \hat{O}, \hat{H}_0 ]$ \\
\hline
\end{tabular}

注意
\begin{itemize}
\item{非可观测量算符不一定遵循上述变化规律，例如密度矩阵$|\psi\rangle\langle\psi|$需要按态的演化方程去计算，而不遵循可观测量的算符演化方程。}
\item{Schr\"odinger绘景和Heisenberg绘景都是Interaction绘景的特例，前者是取了$\hat{H}_0 = 0$，后者是取了$\hat{H}_I = 0$}
\end{itemize}

\ech
\end{frame}

\begin{frame} 
\chtitle{证明一切绘景等价} 
\bch
{\small 
我们只要证明一切Interaction绘景和Sch\"odinger绘景等价，Heisenberg绘景作为Interaction绘景的一种特殊情况则无须再额外证明。所谓两个绘景等价是指任何可观测量的矩阵元$\langle \psi_1 | \hat{O}|\psi_2\rangle$在两个绘景里都相同。
\skipline

把Sch\"odinger绘景下的可观测算符记为$(\hat{O})_S$， 态记为$|\psi\rangle_S$。Interaction绘景下的可观测算符记为$(\hat{O})_I$, 态记为$|\psi\rangle_I$。假设在$t=0$时刻两个绘景下的可观测算符和态均相同，之后分别按各自绘景下的演化方程进行演化。
\begin{itemize}
\item{证明$(\hatH_0)_I$不随时间变化，从而有$(\hatH_0)_I=(\hatH_0)_S$。 对$H_0$我们无须再注明是哪个绘景。}
\item{证明$(\hat{O})_I = e^{\ii\hatH_0 t}(\hat{O})_S e^{-\ii\hatH_0 t}$, 特别地$(\hatH_I)_I= e^{\ii\hatH_0 t}(\hatH_I)_S e^{-\ii\hatH_0 t}$.}
\item{证明$|\psi\rangle_I = e^{\ii\hatH_0 t} |\psi\rangle_S $}
\item{证明对任何态$|\psi_1\rangle$和$|\psi_2\rangle$，$\langle \psi_1 | \hat{O}|\psi_2\rangle$在两个绘景下相同。}
\end{itemize}
}
\ech
\end{frame}

\begin{frame} 
\chtitle{有自相互作用的标量场} 
\bch
考虑拉氏密度为
$$\lagr = \frac{1}{2}\partial^\mu\phi\partial_\mu\phi - \frac{1}{2}m^2\phi^2 - \frac{\lambda}{4!}\phi^4$$
的实标量场。
 
\ech
\end{frame}

\begin{frame} 
\chtitle{} 
\bch

\ech
\end{frame}

\begin{frame} 
\chtitle{} 
\bch

\ech
\end{frame}

\end{document}


